\documentclass[f,proposal,palatino]{AIGpaper}
% Please read the README.md file for additional information on the parameters and overall usage of AIGthesis

\usepackage[english,ngerman]{babel}		    % English and new German spelling
\usepackage[utf8]{inputenc}                 % correct input encoding
\usepackage[T1]{fontenc}                    % correct output encoding
\usepackage{graphicx}					    % enhanced support for graphics
\usepackage{tabularx}				      	% more flexible tabular
\usepackage{amsfonts}					    % math fonts
\usepackage{amssymb}					    % math symbols
\usepackage{amsmath}					    % overall enhancements to math environment
\usepackage{cancel}
\usepackage{soul}
\usepackage{enumitem}
\usepackage{forest}
\usetikzlibrary{tikzmark}
\englishtrue

% optional packages
\usepackage{tikz}                           % creating graphs and other structures
\usetikzlibrary{arrows,positioning}
\tikzset{
    %Define standard arrow tip
    >=stealth',
    %Define style for argument
    args/.style={circle, minimum size=0.9cm,draw=black, thick,fill=white},
}

\newtheorem{definition}{Definition}

% === set information ===
\author{Christoph Kaplan}
        
\title{An Investigation of a Notion of Variable Forgetting that Minimizes Truth Values}


% ========== abstract ===========

\abstract{
This proposal outlines the motivation and problem space of a bachelor thesis focused on exploring a skeptical notion of variable forgetting. We introduce an alternative forget operation and take initial steps towards its investigation, providing an overview of our interests and objectives.
}


\begin{document}
\maketitle % prints title and author information, as well as the abstract 

% optional: change document language from ngerman to english
\selectlanguage{english}


% ========== beginning of the actual text section ===========

\section{Problem and Motivation}
Propositional logic deals, among other things, with the satisfiability of sentences or the inference of sentences from a set of premises. A knowledge base may contain many different statements in various respects. To assert or entail sentences about specific phenomena or interest, it would be advantageous to work only with relevant knowledge to reduce computational complexity. In other words, non-relevant or irrelevant information in relation to the current interest could be "forgotten". Furthermore, knowledge bases change over time, and old knowledge might need to be contracted. To this end, an operation exists for eliminating variables within a proposition. As stated in \cite{lin1994forget} or \cite{lang2003propositional}, we define variable forgetting inductively as follows:

\begin{definition}[Variable Forgetting]
Let F be a formula, p an atomic variable and S as Set of Variables. The \emph{forgetting of p in F $\mathit{Forget}(F,p)$}, respectively, the \emph{forgetting of S in F} $Forget(F,S)$, is given inductively by:
\end{definition}

\begin{itemize}
  \item $\mathit{Forget(F,\emptyset)} = F$ 
  \item $\mathit{Forget(F,\{p\})} = F[p/\top] \lor F[p/\bot]$
  \item $\mathit{Forget}(F,p) = \mathit{Forget}(F,\{p\})$
  \item $\mathit{Forget(F,S \cup \{p\})} = \mathit{Forget(Forget(F,\{p\}), S)}$
\end{itemize}
For some formulas $F$, this operation of variable elimination can lead to a notable reductive or drastic results. We start by considering an example where $\mathit{Forget}$ behaves quite appropriate. Let a sentence $F_1$ be $a \land b$ and we choose to forget about $b$:
\begin{align*}
\mathit{Forget(F_1,b)} &= F_1[b/\top] \lor F_1[b/\bot] \\ 
&\equiv (a \land \top) \lor (a \land \bot) \\ 
&\equiv a \lor \bot \\ 
&\equiv a
\end{align*}
We can see, as shown above, that $\mathit{Forget(F_1,b)}$ is equivalent to $a$ which seems to be the intuitive result when thinking of forgetting about $b$ in $F_1$. Next, we consider forgetting of $b$ in the formula $F_2 = a \lor b$
\begin{align*}
    \mathit{Forget(F_2,b)} &= F_2[b/\top] \lor F_2[b/\bot] \\ 
    &\equiv (a \lor \top) \lor (a \lor \bot) \\ 
    &\equiv \top \lor a \\ 
    &\equiv \top 
\end{align*}
We observe that $\mathit{Forget(F_2,b)}$ is \emph{tautological}, which can be seen as a rather drastic result when we choose to forget $b$ only.

A variation of a forget operation could be defined as follows:
    \begin{definition}[Skeptical Variable Forgetting]
    Let F be a formula, p an atomic variable and S as Set of Variables. The \emph{skeptical forgetting of p in F $\mathit{SkepForget}(F,p)$}, respectively, the \emph{skeptical forgetting of S in F} $\mathit{SkepForget(F,S)}$, is given inductively by:
    \end{definition}
        
    \begin{itemize}
      \item $\mathit{SkepForget(F,\emptyset)} = F$
      \item $\mathit{SkepForget(F,\{p\})} = F[p/\top] \land F[p/\bot]$
        \item $\mathit{SkepForget}(F,p) = \mathit{SkepForget}(F,\{p\})$
      \item $\mathit{SkepForget(F,S \cup \{p\})} = \mathit{SkepForget(SkepForget(F,\{p\}), S)}$
    \end{itemize}

The idea is to use conjunction instead of disjunction. The consequences for our previous examples are the following:

    \begin{align*}
    \mathit{SkepForget(F_1,b)} &= F_1[b/\top] \land F_1[b/\bot] \\ 
    &\equiv (a \land \top) \land (a \land \bot) \\ 
    &\equiv a \land \bot \\ 
    &\equiv \bot
    \end{align*}
    
    \begin{align*}
        \mathit{SkepForget(F_2,b)} &= F_2[b/\top] \land F_2[b/\bot] \\ 
        &\equiv (a \lor \top) \land (a \lor \bot) \\ 
        &\equiv \top \land a \\ 
        &\equiv a 
    \end{align*}
We observe that our alternative operation, at least in these examples, doesn't necessarily entail less drastic results. We obtain a contradiction in the first case and in the second case, our equivalent to $a$. It suggests a duality between the two operations.

Figure \ref{fig:tt1} provides the truth tables for forgetting of b in F and strong forgetting of b in F.
\begin{figure}[ht]
    \centering
       \begin{displaymath}
            \begin{array}{|c c|c|c|c|c|c|c|}
            a & b & F_1 & \mathit{Forget(F_1,b)} & \mathit{SkepForget(F_1,b)} & F_2 & \mathit{Forget(F_2,b)} & \mathit{SkepForget(F_2,b)}\\ 
            \hline
            0 & 0 & 0 & 0 & 0 & 0 & 1 & 0\\
            0 & 1 & 0 & 0 & 0 & 1 & 1 & 0\\
            1 & 0 & 0 & 1 & 0 & 1 & 1 & 1\\
            1 & 1 & 1 & 1 & 0 & 1 & 1 & 1\\
            \end{array}
       \end{displaymath}
    \caption{truth table}
    \label{fig:tt1}
\end{figure}

\section{Epistemic Interest and Questions}
Reflecting upon the subject matter raises the question about the conceptual "essence" of forgetting, and more precisely, the normative expectations entailed by such an operation and its alignment with our definitions. 
Its trivial that a universal comprehension can not be presumed. Since such a discussion might exceed the scope, it would still be helpful, or at least an open question, if we could develop a working thesis to guide our exploration. 
Furthermore, the question of "focus" arises. Our operations not only eliminate variables but also their associated connectives. Intuitively, we do not assign any significance to what this entails. However, it seems to be important to us. Perhaps we can incorporate such considerations into the thesis. Nevertheless, we would like to discuss somewhat more tangible problems as discussed in the following subsections. 

\subsection{Syntactic Questions}
Upon closer examination of more complex formulas, such as $F_3 = (a \lor b) \land c$, a deeper understanding of the operations emerges. We compare the results of forgetting and strong forgetting in $F_3$ to shed light on their effects:

\begin{enumerate}[label=(\arabic*)]
    \item $\mathit{Forget(F_3,a)} \equiv c$
    \item $F\mathit{Forget(F_3,b)} \equiv c$
    \item $\mathit{Forget(F_3,c)} \equiv a \lor b$
    \item $\mathit{SkepForget(F_3,a)} \equiv b \land c$
    \item $\mathit{SkepForget(F_3,b)} \equiv a \land c$
    \item $\mathit{SkepForget(F_3,c)} \equiv \bot$
\end{enumerate}

\noindent If $\top$ and $\bot$ are considered equally "extreme", one could argue that the previous examples were equally drastic in their results. However, the case of $F_3$ reveals a more nuanced behavior. Examples (1) and (2) reduce the formula to just $c$, forgetting two variables instead of one, while example (3) yields an intuitively plausible result. On the other hand, examples (4) and (5) seem to be less reductive and more in line with the concept of forgetting when compared to $\mathit{Forget}$. Nevertheless, example (6) presents a contradiction, which is more drastic than what was observed in examples (1) and (2).

In summary, the $\mathit{SkepForget}$ operation, particularly in the case of $F_3$,  provides us with two desired results, whereas the $\mathit{Forget}$ operation yields only one. However, the singular drastic result produced by the $\mathit{SkepForget}$ operation is more extreme than the drastic result of the $\mathit{Forget}$ operation. 
Furthermore, it lacks the suggested duality observed in earlier instances, thereby prompting further questions for this thesis proposal. Does the outcome predominantly correlate with the top layer of connectives in the syntax tree? What is the effect or influence of deeper nested expressions on the final result? Perhaps the examination of the syntax trees could be enlightening?

\begin{figure}[ht]
\begin{displaymath}
   \begin{array}{c c c c}
   \mathit{Forget(F_3, c)} & \mathit{SkepForget(F_3, c)} & \mathit{Forget(F_3, b)} & \mathit{SkepForget(F_3, b)}\\ 
        \begin{forest}
            [ $\lor$
                [$\land$
                    [$\lor$
                        [$a$]
                        [$b$]
                    ]
                    [$\top$]
                ]
                [$\land$
                    [$\lor$
                        [$a$]
                        [$b$]
                    ]
                    [$\bot$]
                ]
            ]
        \end{forest}         
            & 
        \begin{forest}
            [ $\land$
                [$\land$
                    [$\lor$
                        [$a$]
                        [$b$]
                    ]
                    [$\top$]
                ]
                [$\land$
                    [$\lor$
                        [$a$]
                        [$b$]
                    ]
                    [$\bot$]
                ]
            ]
        \end{forest}
            &
        \begin{forest}
            [ $\lor$
                [$\land$
                    [$\lor$
                        [$a$]
                        [$\top$]
                    ]
                    [$c$]
                ]
                [$\land$
                    [$\lor$
                        [$a$]
                        [$\bot$]
                    ]
                    [$c$]
                ]
            ]
        \end{forest}         
            & 
        \begin{forest}
            [ $\land$
                [$\land$
                    [$\lor$
                        [$a$]
                        [$\top$]
                    ]
                    [$c$]
                ]
                [$\land$
                    [$\lor$
                        [$a$]
                        [$\bot$]
                    ]
                    [$c$]
                ]
            ]
        \end{forest}
        \\
            \equiv c & \equiv a \land c & \equiv a \lor b & \equiv \bot \\     
    \end{array}    
\end{displaymath}   

    \caption{Syntax Trees}
    \label{fig:syntree}
\end{figure}

\noindent Figure \ref{fig:syntree} illustrates that the syntax trees of $\mathit{Forget(F_3,c)}$ exhibit similar syntactic patterns to the previous example of $\mathit{Forget(F_1,b)}$. However, $\mathit{Forget(F_3,b)}$ appears to behave differently upon initial observation, suggesting the need for further investigation to attain a deeper understanding.

\subsection{Semantic Questions}
The truth table presented in Figure \ref{fig:tt1} is not quite appropriate for illustrating the forget operations. Since we essentially can no longer talk about the variable $b$. Therefore, we need to remove $b$ from the table. However, this leads us to another problem concerning semantics. We observe that $b$ is a necessary criteria to distinguish the interpretations. If we remove $b$, we introduce ambiguity in the interpretations as they were dependent on the truth value of $b$. 
Let's visualize and take a closer look at what happens. Figure \ref{fig:ttambiguity} shows a table with the example $F_1$. We remove $b$ and observe that we now have multiple possibilities for $a$, two for the truth value 0, as well as for the value 1.

\begin{figure}[ht]
    \centering
        \begin{displaymath}
            \begin{array}{c||c c|c|c|c|}
            Int & a &\tikzmark{start0} b & F_1 & \mathit{Forget(F_1,b)} & \mathit{SkepForget(F_1,b)}\\ 
            \hline
            1 & 0\tikzmark{arstart1} & 0 &\tikzmark{arend1} 0 & 0 & 0 \\
            2 & 0\tikzmark{arstart2} & 1 &\tikzmark{arend2} 0 & 0 & 0\\
            3 & 1\tikzmark{arstart3} & 0 &\tikzmark{arend3} 0 & 1 & 0 \\
            4 & 1\tikzmark{arstart4} &\tikzmark{end0}1 & \tikzmark{arend4}1 & 1 & 0 \\
            \end{array}
            \tikz[remember picture] \draw[overlay, line width=1pt, red] ([xshift=.25em,yshift=.3em]pic cs:start0) -- ([xshift=.25em]pic cs:end0);
            \tikz[remember picture] \draw[->,overlay, line width=1pt] ([yshift=.35em]pic cs:arstart1) -- ([yshift=.35em]pic cs:arend1);
            \tikz[remember picture] \draw[->,overlay, line width=1pt] ([yshift=.35em]pic cs:arstart2) -- ([yshift=.35em]pic cs:arend2);
            \tikz[remember picture] \draw[->,overlay, line width=1pt] ([yshift=.35em]pic cs:arstart3) -- ([yshift=.35em]pic cs:arend3);
            \tikz[remember picture] \draw[->,overlay, line width=1pt] ([yshift=.35em]pic cs:arstart4) -- ([yshift=.35em]pic cs:arend4);
        \end{displaymath}
    \caption{Ambiguous Assignments}
    \label{fig:ttambiguity}
\end{figure}
    

\noindent To address the semantic issue of ambiguity, we can consider removing or forgetting certain interpretations. However, the question arises: which interpretations should be retained/discarded? To retain universality it's evident that we need to retain one interpretation where $a$ equals 0 and another where $a$ equals 1. 
In our case, lets focus on the possibility that $a$ equals 0, and $b$ is no longer relevant, it leaves us with the decision to choose between 0 or 0 for $F_1$. And for the case $a$ equals 1 we have to decide between 0 or 1 for $F_1$. Figure \ref{fig:poss} shows available options:

\begin{figure}[ht]
        Option 1:
        \begin{displaymath}
            \begin{array}{c||c c|c|c|c|}
            Int & a &\tikzmark{start4} b & F_1 & \mathit{Forget(F_1,b)} & \mathit{SkepForget(F_1,b)}\\ 
            \hline
            1 & 0 & 0 & 0 & 0 & 0 \\
            2 & \tikzmark{start5}0 & 1 & 0 & 0 & 0\tikzmark{end5}\\
            3 & 1 & 0 & 0 & 1 & 0 \\
            4 & \tikzmark{start6}1 &\tikzmark{end4}1 & 1 & 1 & 0\tikzmark{end6} \\
            \end{array}
            \tikz[remember picture] \draw[overlay, line width=1pt, red] ([xshift=.25em,yshift=.3em]pic cs:start4) -- ([xshift=.25em]pic cs:end4);
            \tikz[remember picture] \draw[overlay, line width=1pt, red] ([yshift=.35em]pic cs:start5) -- ([yshift=.35em]pic cs:end5);
            \tikz[remember picture] \draw[overlay, line width=1pt, red] ([yshift=.35em]pic cs:start6) -- ([yshift=.35em]pic cs:end6);
        \end{displaymath}

        Option 2:
        
        \begin{displaymath}
            \begin{array}{c||c c|c|c|c|}
            Int & a &\tikzmark{start1} b & F_1 & \mathit{Forget(F_1,b)} & \mathit{SkepForget(F_1,b)}\\ 
            \hline
            1 & \tikzmark{start2}0 & 0 & 0 & 0 & 0 \tikzmark{end2}\\
            2 & 0 & 1 & 0 & 0 & 0\\
            3 & \tikzmark{start3}1 & 0 & 0 & 1 & 0 \tikzmark{end3}\\
            4 & 1 &\tikzmark{end1}1 & 1 & 1 & 0 \\
            \end{array}
            \tikz[remember picture] \draw[overlay, line width=1pt, red] ([xshift=.25em,yshift=.3em]pic cs:start1) -- ([xshift=.25em]pic cs:end1);
            \tikz[remember picture] \draw[overlay, line width=1pt, red] ([yshift=.35em]pic cs:start2) -- ([yshift=.35em]pic cs:end2);
            \tikz[remember picture] \draw[overlay, line width=1pt, red] ([yshift=.35em]pic cs:start3) -- ([yshift=.35em]pic cs:end3);
        \end{displaymath}
    \caption{Collapse Options}
    \label{fig:poss}
\end{figure}

\noindent For the first option ($b$ was 0), we choose to preserve interpretations 1 and 3. $F_1$ was minimal with respect to $a$. For the second option (if $b$ was 1). We choose the interpretations 2 and 4 to be retained, where $F_1$ was maximal in respect to $a$. The question is weather there is one correct option or are they contingent?\\

In \cite{lang2003propositional}, the authors denote formal properties of the $\mathit{Forget}$ operation, including: $Mod(Forget(F, x)) = Mod(F) \cup \{Switch(w, x) | w \models F\}$, where $Switch(w, x)$ denotes the interpretation that maintains the same truth value as $w$ for all variables except $x$, and assigns to $x$ the opposite truth value given by $w$. This function, $Switch(w, x)$, appears to offer a promising approach to resolving the ambiguity discussed earlier.
We aim to uncover analogous formal properties for our variation as well. For instance $Mod(SkepForget(F, x)) = Mod(F) \cap \{Switch(w, x) | w \models F\}$ seems to be a good candidate.

\section{Course of action}
Our timeline begins with the initiation of thesis work on 1. April 2024. We plan to introduce the thesis proposal through a 15-minute presentation at the Oberseminar on [date]. Throughout the thesis development process, we will maintain regular communication with our supervisor until [date]. As the project nears its conclusion, we will proceed to finalize the thesis with minimal supervision and submit it online. Finally, we will present the work by a 30-minute presentation of the thesis at the Oberseminar on [date].


% References
\bibliographystyle{alpha}
\bibliography{references}

\end{document}
